\documentclass[10pt,a4paper]{amsart}
\usepackage{amsmath,amssymb,amsthm,graphicx}
% Retained minimal preamble for the Stacks Project chapter sites-cohomology.tex (Cohomology on Sites).
% Verbatim excerpt of the chapter preamble: Xy-pic setup (native preamble lines 14-18) and the native
% theorem declarations with the mathematical macros (native preamble lines 154-217). Copied mechanically;
% nothing added, removed or redefined. The wrapper class and the amsmath/amssymb/amsthm/graphicx base
% packages are supplied by the assembler, so the class/externaldocument part of the native preamble is omitted.
% For commutative diagrams we use Xy-pic
\usepackage[all]{xy}

% We use 2cell for 2-commutative diagrams.
\xyoption{2cell}
% Theorem environments.
%
\theoremstyle{plain}
\newtheorem{theorem}[subsection]{Theorem}
\newtheorem{proposition}[subsection]{Proposition}
\newtheorem{lemma}[subsection]{Lemma}

\theoremstyle{definition}
\newtheorem{definition}[subsection]{Definition}
\newtheorem{example}[subsection]{Example}
\newtheorem{exercise}[subsection]{Exercise}
\newtheorem{situation}[subsection]{Situation}

\theoremstyle{remark}
\newtheorem{remark}[subsection]{Remark}
\newtheorem{remarks}[subsection]{Remarks}

\numberwithin{equation}{subsection}

% Macros
%
\def\lim{\mathop{\mathrm{lim}}\nolimits}
\def\colim{\mathop{\mathrm{colim}}\nolimits}
\def\Spec{\mathop{\mathrm{Spec}}}
\def\Hom{\mathop{\mathrm{Hom}}\nolimits}
\def\Ext{\mathop{\mathrm{Ext}}\nolimits}
\def\SheafHom{\mathop{\mathcal{H}\!\mathit{om}}\nolimits}
\def\SheafExt{\mathop{\mathcal{E}\!\mathit{xt}}\nolimits}
\def\Sch{\mathit{Sch}}
\def\Mor{\mathop{\mathrm{Mor}}\nolimits}
\def\Ob{\mathop{\mathrm{Ob}}\nolimits}
\def\Sh{\mathop{\mathit{Sh}}\nolimits}
\def\NL{\mathop{N\!L}\nolimits}
\def\CH{\mathop{\mathrm{CH}}\nolimits}
\def\proetale{{pro\text{-}\acute{e}tale}}
\def\etale{{\acute{e}tale}}
\def\QCoh{\mathit{QCoh}}
\def\Ker{\mathop{\mathrm{Ker}}}
\def\Im{\mathop{\mathrm{Im}}}
\def\Coker{\mathop{\mathrm{Coker}}}
\def\Coim{\mathop{\mathrm{Coim}}}

% Boxtimes
%
\DeclareMathSymbol{\boxtimes}{\mathbin}{AMSa}{"02}

%
% Macros for moduli stacks/spaces
%
\def\QCohstack{\mathcal{QC}\!\mathit{oh}}
\def\Cohstack{\mathcal{C}\!\mathit{oh}}
\def\Spacesstack{\mathcal{S}\!\mathit{paces}}
\def\Quotfunctor{\mathrm{Quot}}
\def\Hilbfunctor{\mathrm{Hilb}}
\def\Curvesstack{\mathcal{C}\!\mathit{urves}}
\def\Polarizedstack{\mathcal{P}\!\mathit{olarized}}
\def\Complexesstack{\mathcal{C}\!\mathit{omplexes}}
% \Pic is the operator that assigns to X its picard group, usage \Pic(X)
% \Picardstack_{X/B} denotes the Picard stack of X over B
% \Picardfunctor_{X/B} denotes the Picard functor of X over B
\def\Pic{\mathop{\mathrm{Pic}}\nolimits}
\def\Picardstack{\mathcal{P}\!\mathit{ic}}
\def\Picardfunctor{\mathrm{Pic}}
\def\Deformationcategory{\mathcal{D}\!\mathit{ef}}
\title{The qc coverings on the category of locally quasi-compact Hausdorff spaces form a topology}
\author{The Stacks Project Authors}
\date{stacks-project:sites-cohomology.tex:local-snapshot-20261011; GNU FDL 1.2 or later}
\begin{document}
\maketitle
\noindent\textit{Editorial note.} All delivered mathematical text is the author's own, copied from the pinned Stacks Project chapter file sites-cohomology.tex; no statement or proof step was rewritten, reordered or removed.\par\medskip
\noindent\textit{Editorial note.} The proof is the author's complete argument for the three parts, including the quasi-compactness reduction inside part (2) that the author writes out in full.\par\medskip
\noindent\textit{Editorial note.} The label lemma-qc occurs in no other staged Stacks chapter, so no cross-chapter restatement had to be weighed for this card. In particular there is no sibling statement of the same name in the chapter etale-cohomology.tex.\par\medskip
\noindent\textit{Editorial note.} The preamble is the chapter's own retained minimal preamble (sources/stacks-sites-cohomology/retained-preamble.tex), a verbatim excerpt of the native preamble: the Xy-pic setup and the native theorem declarations with the mathematical macros.\par\medskip
\noindent\textit{Allowed standard prerequisites (editorial).}\begin{itemize}
\item The category LC of Hausdorff and locally quasi-compact spaces. Same chapter, the definition of the category LC with morphisms being continuous maps. Supplies the ambient category on which the qc coverings are declared.
\item qc covering. Same chapter, the definition of a qc covering: a family \{f\_i : X\_i -> X\} such that every point of X has a neighbourhood covered by finitely many quasi-compact subsets of the X\_i. Supplies the notion verified in all three parts, and the exact shape of the neighbourhoods used in parts (2) and (3).
\item Quasi-compactness is stable under fibre products in LC. Same chapter, the permanence of quasi-compactness used when the base change F x\_X E\_j is called quasi-compact. Supplies the quasi-compactness of F x\_X E\_j inside X' x\_X X\_\{i\_j\} in part (3).
\end{itemize}\medskip
\begin{lemma}
\label{lemma-qc}
Let $X$ be a Hausdorff and locally quasi-compact space, in other words,
an object of $\textit{LC}$.
\begin{enumerate}
\item If $X' \to X$ is an isomorphism in $\textit{LC}$ then
$\{X' \to X\}$ is a qc covering.
\item If $\{f_i : X_i \to X\}_{i\in I}$ is a qc covering and for each
$i$ we have a qc covering $\{g_{ij} : X_{ij} \to X_i\}_{j\in J_i}$, then
$\{X_{ij} \to X\}_{i \in I, j\in J_i}$ is a qc covering.
\item If $\{X_i \to X\}_{i\in I}$ is a qc covering
and $X' \to X$ is a morphism of $\textit{LC}$ then
$\{X' \times_X X_i \to X'\}_{i\in I}$ is a qc covering.
\end{enumerate}
\end{lemma}

\begin{proof}
Part (1) holds by the remark above that open coverings are qc coverings.

\medskip\noindent
Proof of (2). Let $x \in X$. Choose $i_1, \ldots, i_n \in I$ and
$E_a \subset X_{i_a}$ quasi-compact such that $\bigcup f_{i_a}(E_a)$
is a neighbourhood of $x$. For every $e \in E_a$ we can find
a finite subset $J_e \subset J_{i_a}$ and quasi-compact
$F_{e, j} \subset X_{ij}$, $j \in J_e$ such that $\bigcup g_{ij}(F_{e, j})$
is a neighbourhood of $e$. Since $E_a$ is quasi-compact we find
a finite collection $e_1, \ldots, e_{m_a}$ such that
$$
E_a \subset
\bigcup\nolimits_{k = 1, \ldots, m_a}
\bigcup\nolimits_{j \in J_{e_k}} g_{ij}(F_{e_k, j})
$$
Then we find that
$$
\bigcup\nolimits_{a = 1, \ldots, n}
\bigcup\nolimits_{k = 1, \ldots, m_a}
\bigcup\nolimits_{j \in J_{e_k}} f_i(g_{ij}(F_{e_k, j}))
$$
is a neighbourhood of $x$.

\medskip\noindent
Proof of (3). Let $x' \in X'$ be a point. Let $x \in X$ be its image.
Choose $i_1, \ldots, i_n \in I$ and quasi-compact subsets
$E_j \subset X_{i_j}$ such that $\bigcup f_{i_j}(E_j)$ is a
neighbourhood of $x$. Choose a quasi-compact neighbourhood $F \subset X'$
of $x'$ which maps into the quasi-compact neighbourhood
$\bigcup f_{i_j}(E_j)$ of $x$. Then
$F \times_X E_j \subset X' \times_X X_{i_j}$ is a
quasi-compact subset and $F$ is the image of the map
$\coprod F \times_X E_j \to F$. Hence the base change is a
qc covering and the proof is finished.
\end{proof}

\end{document}
